\documentclass[FCD_GNN.tex]{subfiles}

\begin{document}
\chapter{Discussion}
\label{chapter:Discussion}

The experiments conducted in this thesis demonstrate that incorporating textual information into surface-based GNN architectures can substantially improve the detection of FCD type II. At the same time, the results highlight important trade-offs that need to be considered when designing multimodal models.

A more detailed analysis of the GNN experiments reveals a nuanced picture. 
Integrating additional MELD feature stages into the GNN blocks increases the model’s ability to detect lesions:
sensitivity and lesion-level coverage improved as more stages were incorporated, indicating that neighborhood aggregation 
can amplify subtle abnormal patterns that may be missed by purely local representations. 
This benefit, however, comes at the cost of reduced Specificity, i.e., a higher false-positive rate on healthy controls.
This is an important limitation for clinical use, since our evaluation does not include other epilepsy-related abnormalities or non-FCD lesion types; therefore, we cannot assess how often the model would spuriously highlight regions in patients with different pathologies.

Importantly, the results show that adding more GNN layers does not monotonically improve performance.
The effectiveness of message passing strongly depends on the structure of the underlying MELD feature stages.
For higher-dimensional and sparser stages, neighborhood aggregation stabilizes features and enriches contextual information,
making GNN layers beneficial.
In contrast, for lower-dimensional stages (e.g., stages 6--7), where the graph becomes denser and representations are more compressed,
additional propagation leads to oversmoothing, blurring spatial distinctions and ultimately degrading performance.


In the warm--up experiments, we observed that freezing the visual backbone for the first few epochs can be beneficial. 
A plausible explanation is that the model becomes less sensitive to prompt variability and even mildly incorrect prompts. 
This behavior can be attributed to several factors:
\begin{enumerate}
  \item \textbf{Partial but useful signal.} Even an inexact prompt still conveys coarse structural information (e.g., hemisphere or lobe). 
  Such prompts can narrow the search space and reduce the need to attend over the entire cortex.

  \item \textbf{Decoder robustness.} Owing to pretraining, the decoder can tolerate mild text--image mismatch. 
  When the prompt is uninformative, the model relies more on visual features and treats the text mainly as a weak attention prior.

  \item \textbf{Regularization effect.} Training with slightly noisy prompts discourages the model from overfitting to specific wording and promotes more general text--image associations. 
  This acts as mild regularization: predictions become more conservative, reducing false positives, which can improve precision (e.g., PPV), while sensitivity often remain stable.
\end{enumerate}

Experiments with different forms of textual input showed that even coarse prompts (hemisphere- or lobe-only) can meaningfully improve lesion detection compared to purely vision-based baselines.
Across configurations, gains in sensitivity typically come with a trade-off in specificity, i.e., an increased number of false positive clusters.
Overall, reduced-text prompts (e.g., hemisphere + lobe) often provide a more favorable balance between detecting lesions and limiting false positives, which may be preferable for clinical use.

Several limitations must be acknowledged. 
While the MELD surface features are ComBat-harmonized to reduce inter-site effects, residual variability may remain due to differences in acquisition protocols, preprocessing, and lesion annotation practices across centers.
Furthermore, the MELD cohort predominantly consists of histologically confirmed type II FCD cases, and it therefore remains unclear whether the same conclusions extend to other FCD subtypes.
Finally, the model relies on atlas-based text generated automatically; incorporating free-form radiology reports could yield a more realistic estimate of performance under clinical conditions, where textual descriptions are not always precise.

In this study, we did not explore the effect of unfreezing different numbers of LLM layers or varying the number of cross-modal connections to the GuideDecoder. However, previous work (Huemann et al., 2024 \cite{Huemann2024ConTEXTualNet}) shows that selectively unfreezing specific layers of a language model and fine-tuning them for a segmentation task can yield consistent performance gains. 
Consequently, systematic experiments on partial LLM fine-tuning, on alternative strategies for connecting the text encoder to the GuideDecoder, and on increasing the architectural capacity of the GuideDecoder itself represent promising directions for future research.

Another conceptual direction concerns the representation of anatomical descriptions. 
In this thesis, atlas information was encoded via textual prompts and processed by a pretrained language model. 
An alternative would be to express the same information as a numerical feature vector with one dimension per cortical region. 
In such a vector, the percentages mentioned in the textual description (e.g., 52.34\% Right Frontal Orbital Cortex; 26.06\% Right Frontal Pole) would directly populate the corresponding entries, while all other regions would be set to zero. 
For coarse labels such as lobe-only or hemisphere-only descriptions, the corresponding weight could be uniformly distributed across all constituent regions.

While this vector--based formulation is simple, interpretable, and avoids the computational overhead of a pretrained LLM, it is also less expressive in its basic form. 
When used without an explicitly defined structure, a model trained solely on such numerical representations treats each vector dimension as an independent feature and does not explicitly encode anatomical relationships, such as lobar or hemispheric grouping, spatial adjacency, or hierarchical organization. 
Although these relationships could in principle be incorporated through carefully designed numerical encodings or structured representations, this requires additional modeling assumptions and design choices.
In contrast, language models implicitly encode semantic relationships between anatomical terms, including regional similarity and hierarchical structure, and naturally handle coarse, uncertain, or partially specified descriptions. 
This makes LLM--based representations particularly suitable for prototyping multimodal fusion and exploring how textual guidance can complement visual features. 
Accordingly, while numerical vector representations remain a sustainable baseline, their relative performance compared to LLM-based approaches depends on how explicitly anatomical structure and uncertainty are modeled, and remains an open question for future work.

Overall, the discussion of results highlights that textual guidance is a promising direction for FCD detection. 
However, the balance between sensitivity and precision remains a key design choice that must be adapted depending on whether the clinical task prioritizes coverage or specificity.

% \item How does fine-tuning different numbers of layers in a pre-trained \ac{llm} affect model performance?
% \item Performed extensive experiments to evaluate different design choices, such as the number of connections in the \ac{gnn} blocks and the use of various textual branches in the GuideDecoder, to improve segmentation results.
% Chapter~\ref{chapter:Experiments} details the experimental setup (losses, metrics, and implementation specifics) and reports the main results: baseline comparisons against MELD, the effect of linking different numbers of MELD feature stages to the \acs{gnn} block, the influence of (partially) unfreezing RadBERT and using text guidance in the decoder, and ablations on the number of text–GuideDecoder connections; where relevant, we also compare pre- and post-fine-tuning performance.
\end{document}
